\par\addvspace{12pt}\Needspace{9\baselineskip}

\section{圆锥曲线}

% ITEM: quadratic-example-021
\solutionitem{quadratic-example-021}{例4.21}
圆锥面为\(X^2+Y^2=Z^2\)。分别求截平面\(Z=2\)、\(Z=X/2+2\)、\(Z=X+2\)、\(Z=2X+2\)上的截线方程，判断类型，并解释临界倾角。
\solutionanswer
逐一将平面关系代入锥面，得到
\[
\begin{array}{c|c|c}
\text{平面}&\text{截线的投影方程}&\text{类型}\\\hline
Z=2&X^2+Y^2=4&\text{圆}\\
Z=X/2+2&\frac34(X-4/3)^2+Y^2=16/3&\text{椭圆}\\
Z=X+2&Y^2=4(X+1)&\text{抛物线}\\
Z=2X+2&3(X+4/3)^2-Y^2=4/3&\text{双曲线}
\end{array}
\]
各方程均须连同相应截平面理解。锥母线与轴成\(45^\circ\)，当截平面倾角也为\(45^\circ\)时，平面平行于一条母线，二次平方项恰有一个消失，这是临界情形。


% ITEM: quadratic-thought-03
\solutionitem{quadratic-thought-03}{思考：截平面的退化情形}
不经过锥顶的截平面能否截出两条平行直线？为什么仅改变倾角，而不说明是否经过锥顶，还不能完整确定截线类型？
\solutionanswer
不能。锥面上包含的直线是经过锥顶的母线；若截线含两条完整平行直线，截平面就不得不经过锥顶，与题设矛盾。一般地，倾角决定与母线相对的方向，是否过锥顶则决定相应类型是否退化，例如同一临界倾角可给出抛物线，也可给出一条重母线。


% ITEM: quadratic-exercise-12-01
\solutionitem{quadratic-exercise-12-01}{习题4.12（1）}
圆锥母线与轴成\(30^\circ\)。截平面分别与水平面成\(0^\circ,45^\circ,60^\circ,75^\circ\)，且不经过锥顶，判断各截线类型。
\par
\solutionanswer
母线与轴夹角\(\alpha=30^\circ\)，故临界截平面倾角为\(90^\circ-\alpha=60^\circ\)。因此倾角\(0^\circ\)时为圆；\(45^\circ\)小于临界角，为非圆椭圆；\(60^\circ\)恰为临界角，为抛物线；\(75^\circ\)大于临界角，为双曲线。题设截平面均不经过锥顶，故均为非退化截线。


% ITEM: quadratic-exercise-12-02
\solutionitem{quadratic-exercise-12-02}{习题4.12（2）}
对圆锥面\(X^2+Y^2=4Z^2\)，讨论截平面\(Z=mX+1\)随\(m\ge0\)变化时的截线类型，再将截平面改为\(Z=mX\)比较。
\par
\solutionanswer
圆锥面参数\(k=2\)，截平面倾角满足\(m=\tan\beta\)，临界值\(km=1\)，即\(m=1/2\)。代入\(Z=mX+1\)得
\[(1-4m^2)X^2+Y^2-8mX-4=0.\]
当\(0\le m<1/2\)，二次项同号且配方后右端为正，截线为椭圆（\(m=0\)时为圆）；当\(m=1/2\)时只剩\(Y^2\)平方项且\(X\)一次项非零，为抛物线；当\(m>1/2\)时二次项异号，配方后常数非零，为双曲线。
若改为过锥顶的平面\(Z=mX\)，方程为\((1-4m^2)X^2+Y^2=0\)。于是\(m<1/2\)时交集只有锥顶；\(m=1/2\)时为重直线\(Y=0\)；\(m>1/2\)时为两条相交母线。比较可见，倾角决定二次部分的符号类型，平面是否过顶点则决定相应截线是否退化。


% ITEM: quadratic-exercise-12-03
\solutionitem{quadratic-exercise-12-03}{习题4.12（3）}
将圆锥面与截平面同时以锥顶为中心放大\(t\)倍（\(t>0\)），说明截线如何变化。对一组相互平行且不经过锥顶的截平面，所得截线之间有什么关系？
\par
\solutionanswer
以锥顶为原点，位似 \(X\mapsto tX\) 保持锥面，把截线送到新截线，各长度乘 \(t\)。若两个平行截平面写成 \(n\cdot X=d_1,d_2\)，取 \(t=d_2/d_1\) 即可将前者映到后者；\(t<0\) 时再计入关于锥顶的中心对称。故两截线总相似，相似比为 \(|d_2/d_1|\)。


% ITEM: quadratic-exercise-12-04
\solutionitem{quadratic-exercise-12-04}{习题4.12（4）}
（据2015年浙江卷文科第7题）斜线段 \(AB\) 与平面 \(\alpha\) 成 \(60^\circ\)，\(B\) 为斜足。平面内的动点 \(P\) 满足 \(\angle PAB=30^\circ\)，其轨迹是什么曲线？
\par
\solutionanswer
先不限制 \(P\) 在平面内。射线 \(AP\) 绕射线 \(AB\) 旋转，扫出半顶角为 \(30^\circ\) 的单叶锥面。平面 \(\alpha\) 在 \(B\) 处与锥轴相交，所成角 \(60^\circ\) 大于半顶角而小于直角，故截得的轨迹是非圆椭圆。


% ITEM: quadratic-example-022
\solutionitem{quadratic-example-022}{例4.22 *}
对\(Q=X^2+Y^2-Z^2\)，分别将它限制到平面\(Z=1\)、\(Z=X+1\)、\(X=1\)上，求二次部分与完整齐次矩阵的秩，并与截线类型对应。
\solutionanswer
依次选平面坐标\((X,Y)\)、\((X,Y)\)、\((Y,Z)\)，限制式为\(u^2+v^2-1\)、\(v^2-2u-1\)、\(u^2-v^2+1\)。二次部分的秩分别为\(2,1,2\)，完整齐次矩阵的秩均为\(3\)；截线依次为圆、抛物线、双曲线。第二个完整矩阵为\(\begin{bmatrix}0&0&-1\\0&1&0\\-1&0&-1\end{bmatrix}\)，说明二次部分退化不意味着整条截线退化。


% ITEM: quadratic-exercise-13-01
\solutionitem{quadratic-exercise-13-01}{习题4.13（1） *}
设\(k=1\)。求\(\beta=30^\circ,45^\circ,60^\circ\)时\(B^THB\)的特征值及正负系数个数。
\par
\solutionanswer
两个特征值为\(\cos2\beta,1\)，依次为\((1/2,1),(0,1),(-1/2,1)\)。正、负系数个数依次为\((2,0),(1,0),(1,1)\)，中间情形另有一个零方向。


% ITEM: quadratic-exercise-13-02
\solutionitem{quadratic-exercise-13-02}{习题4.13（2） *}
设截平面\(Z=X+1\)上的点写成\(P=(u,v,u+1)^T\)。求\(B\)、\(P_0\)及式\textup{(4.64)}，说明这里的\(u\)与单位坐标之间的关系。
\par

\par\smallskip\noindent 正文所引关系为：
\begin{equation*}
Q(P_0+BU)=U^T(B^THB)U+2P_0^THBU+P_0^THP_0 .
\end{equation*}
\solutionanswer
取\(P_0=(0,0,1)^T\)，以\((u,v)\)为参数时，\(B\)的两列分别是\((1,0,1)^T,(0,1,0)^T\)。代入锥面的二次型得
\[
 Q(P)=u^2+v^2-k^2(u+1)^2=(1-k^2)u^2+v^2-2k^2u-k^2.
\]
当\(k=1\)时，上式化为\(v^2-2u-1\)。两列正交，长度分别为\(\sqrt2,1\)；若按正文取单位基，第一列改为\((1,0,1)^T/\sqrt2\)，相应坐标为\((\sqrt2\,u,v)\)。


% ITEM: quadratic-exercise-13-03
\solutionitem{quadratic-exercise-13-03}{习题4.13（3） *}
证明截平面的两个正交单位方向改变后，限制二次型的矩阵虽然改变，其秩与正负系数个数仍不变。
\par
\solutionanswer
另取正交单位基时\(B'=BR\)，其中\(R\)为二阶正交矩阵。于是\(B'^THB'=R^T(B^THB)R\)，两限制矩阵合同，故秩与正负系数个数不变。


% ITEM: quadratic-exercise-13-04
\solutionitem{quadratic-exercise-13-04}{习题4.13（4） *}
在式\textup{(4.66)}中取\(k=1\)，分别讨论\(\beta=0^\circ,45^\circ,90^\circ\)。求\(d\ne0\)时的截线方程及类型，再与\(d=0\)比较。
\par

\par\smallskip\noindent 正文所引关系为：
\begin{equation*}
\lambda u^2+v^2-2d(1+k^2)\sin\beta\cos\beta\,u
    +d^2(\sin^2\beta-k^2\cos^2\beta)=0 .
\end{equation*}
\solutionanswer
\(k=1\)时，限制式为\(\cos2\beta\,(u^2-d^2)+v^2-2d\sin2\beta\,u=0\)。三角度分别给出\(u^2+v^2=d^2\)、\(v^2=2du\)、\(u^2-v^2=d^2\)，\(d\ne0\)时依次为圆、抛物线、双曲线。\(d=0\)时依次成为单点、重直线和相交直线对。


% ITEM: quadratic-exercise-13-05
\solutionitem{quadratic-exercise-13-05}{习题4.13（5） *}
在同一截平面上，将原点从\(P_0\)移到\(P_0+BW\)。证明限制式的二次部分仍为\(B^THB\)，一次项系数列则增加\(2B^THBW\)。联系4.3的中心条件，说明何时能把一次项全部消去。
\par
\solutionanswer
记原限制式为\(U^TKU+l^TU+c\)，其中\(K=B^THB\)。新坐标满足\(U=V+W\)，展开得二次部分仍为\(K\)，一次项系数为\(l+2KW\)。故全部消去一次项当且仅当\(2KW+l=0\)有解；\(K\)可逆时解唯一，不可逆时需要\(-l/2\)属于\(K\)的像。


% ITEM: quadratic-example-023
\solutionitem{quadratic-example-023}{例4.23}
一条椭圆截线对应的两个相切圆沿母线的间距为\(10\)，两焦点距离为\(6\)。求截线上位于两焦点连线上的两个点到各焦点的距离，并说明使用了哪些相切关系。
\solutionanswer
记两球与截平面的切点为 \(F_1,F_2\)，过 \(P\) 的母线与两球的切点为 \(T_1,T_2\)。从 \(P\) 向同一球引出的切线段等长，故 \(PF_i=PT_i\)。又因 \(P\) 在 \(T_1,T_2\) 之间，两个焦半径之和为 \(T_1T_2=10\)。

当 \(P,F_1,F_2\) 共线时，若 \(P\) 在两焦点之间，焦半径之和只能是 \(6\)。所以 \(P\) 在焦点线段的延长线上，较长焦半径与较短焦半径相差 \(6\)。由和为 \(10\)、差为 \(6\)，得两段长为 \(8,2\)；换到另一端点，两焦点对应的距离互换。


% ITEM: quadratic-exercise-14-01
\solutionitem{quadratic-exercise-14-01}{习题4.14（1）}
证明从球外一点引出的任意两条切线段等长，并用球心到切点的半径说明理由。
\par
\solutionanswer
设球心为\(O\)、半径为\(R\)，切点为\(T\)。半径\(OT\)垂直于切线，所以\(|PT|^2=|PO|^2-R^2\)。右边与切点无关，因此任意两条切线段等长。


% ITEM: quadratic-exercise-14-02
\solutionitem{quadratic-exercise-14-02}{习题4.14（2）}
在椭圆情形，设两个相切圆所在水平面的高度差为\(h\)，圆锥半顶角为\(\alpha\)。求\(|PF_1|+|PF_2|\)。在双曲线情形，相同的高度差给出怎样的距离差？
\par
\solutionanswer
相切圆的两个接触点在同一母线上，其沿母线的间距为\(h/\cos\alpha\)。由同一点向球所引切线段相等，椭圆的距离和为\(h/\cos\alpha\)；双曲线的距离差绝对值为同一常数\(h/\cos\alpha\)。


% ITEM: quadratic-exercise-14-03
\solutionitem{quadratic-exercise-14-03}{习题4.14（3）}
在不经过锥顶的非圆椭圆截线中，说明距离和为什么严格大于两焦点距离；在双曲线截线中，距离差的绝对值为什么严格小于两焦点距离？
\par
\solutionanswer
取截线上一个不在焦点连线上的点，三角不等式严格成立，故椭圆的距离和严格大于两焦点距离，双曲线的距离差绝对值严格小于两焦点距离。相应和或差在整条截线上为常数，因此得到对曲线本身的结论。


% ITEM: quadratic-example-024
\solutionitem{quadratic-example-024}{例4.24}
（1）圆锥母线与轴成\(45^\circ\)，截平面不经过锥顶。分别求\(\beta=30^\circ,45^\circ,60^\circ,90^\circ\)时的离心率，并判断曲线类型。

（2）平面上取\(F=(0,0)\)、\(\ell:x=3\)，动点\(P\)满足\(|PF|=\frac23d(P,\ell)\)。不先求曲线的标准方程，求\(|PF|\)的取值范围，指出取得两端值的点。
\solutionanswer
（1）用\(e=\sqrt2\sin\beta\)，依次得\(1/\sqrt2,1,\sqrt{3/2},\sqrt2\)，类型依次为椭圆、抛物线、双曲线、双曲线。

（2）记 \(r=|PF|\)。点 \(P=(x,y)\) 不可能在准线右侧或准线上：若 \(x\ge3\)，则 \(r=\tfrac23(x-3)<x\)，与 \(r\ge|x|\) 矛盾。故 \(x=3-3r/2\)。半径为 \(r\) 的圆上能取到这个横坐标，当且仅当
\[
   -r\le 3-3r/2\le r,\qquad\text{即}\qquad 6/5\le r\le6 .
\]
两端分别对应 \(P=(6/5,0),(-6,0)\)。


% ITEM: quadratic-exercise-15-01
\solutionitem{quadratic-exercise-15-01}{习题4.15（1）}
点\(P\)到焦点与准线的距离分别为\(3\)、\(5\)。求离心率并判断类型。若焦半径变为\(6\)，点到同一准线的距离是多少？
\par
\solutionanswer
离心率为\(e=|PF|/d(P,\ell)=3/5<1\)，故为椭圆。保持焦点、准线和离心率不变时，焦半径与到准线距离之比仍为\(3/5\)；若焦半径为\(6\)，则到准线距离为\(6/(3/5)=10\)。


% ITEM: quadratic-exercise-15-02
\solutionitem{quadratic-exercise-15-02}{习题4.15（2）}
圆锥半顶角为\(60^\circ\)。若截线离心率为\(1\)、\(\sqrt3\)，分别求截平面倾角。该圆锥能否截出离心率为\(3\)的双曲线？
\par
\solutionanswer
本节关系为\(e=\sin\beta/\cos\alpha\)。当\(\alpha=60^\circ\)时，\(\cos\alpha=1/2\)，故\(\sin\beta=e/2\)。\(e=1\)时\(\beta=30^\circ\)；\(e=\sqrt3\)时\(\beta=60^\circ\)。而\(0\le\sin\beta\le1\)，故该圆锥能实现的离心率至多为\(2\)，不能截出离心率为\(3\)的双曲线。


% ITEM: quadratic-exercise-15-03
\solutionitem{quadratic-exercise-15-03}{习题4.15（3）}
一条非圆椭圆和一条双曲线分别满足\(|PF|=e\,d(P,\ell)\)。比较曲线上点到焦点、准线的距离大小，并说明这一判断是否取决于点在曲线上的位置。
\par
\solutionanswer
椭圆离心率满足\(0<e<1\)，因此对曲线上每一点均有\(|PF|=e\,d(P,\ell)<d(P,\ell)\)。双曲线满足\(e>1\)，故每一点均有\(|PF|=e\,d(P,\ell)>d(P,\ell)\)。比较只由曲线的固定离心率决定，不依赖点在曲线上的位置。


% ITEM: quadratic-exercise-15-04
\solutionitem{quadratic-exercise-15-04}{习题4.15（4）}
设\(F\)到准线\(\ell\)的距离为\(d\)，曲线满足\(|PF|=e\,d(P,\ell)\)，其中\(0<e<1\)。设\(P\)在\(\ell\)上的垂足为\(Q\)，比较\(|FQ|\)、\(|PQ|\)与\(|PF|\)，求\(|PF|\)的最小可能值，并说明这个值怎样取得。
\par
\solutionanswer
记\(r=|PF|,s=|PQ|\)，则\(r=es<s\)，且三角不等式给出\(s-r\le|FQ|\le s+r\)。又\(|FQ|\ge d\)，所以\(d\le r+r/e\)，即\(r\ge ed/(1+e)\)。设\(H\)是\(F\)在准线上的垂足，取\(P\)在线段\(FH\)上且\(|FP|=ed/(1+e),|PH|=d/(1+e)\)，便满足焦准关系并取到最小值。


% ITEM: quadratic-example-025
\solutionitem{quadratic-example-025}{例4.25}
以\(F(0,0)\)为焦点、\(x=-2\)为准线，分别取\(e=1/2,1,2\)。求三条曲线的方程，化为标准形式，并求其中心或对称轴。
\solutionanswer
统一从\(x^2+y^2=e^2(x+2)^2\)出发。\(e=1/2\)时为
\[
\frac{(x-2/3)^2}{16/9}+\frac{y^2}{4/3}=1,
\]
中心\((2/3,0)\)。\(e=1\)时为\(y^2=4(x+1)\)，对称轴\(y=0\)。\(e=2\)时为
\[
\frac{(x+8/3)^2}{16/9}-\frac{y^2}{16/3}=1,
\]
中心\((-8/3,0)\)，两中心曲线的坐标方向也给出它们的对称轴。


% ITEM: quadratic-exercise-16-01
\solutionitem{quadratic-exercise-16-01}{习题4.16（1）}
求到\(F(1,0)\)的距离等于到直线\(x=-1\)距离的点的轨迹方程，并指出对称轴。
\par
\solutionanswer
焦点为\(F(1,0)\)，准线为\(x=-1\)，离心率\(e=1\)。焦准线定义给出
\[
\sqrt{(x-1)^2+y^2}=|x+1|.
\]
两边非负，平方等价，化简得\(y^2=4x\)。轨迹是以\(x\)轴为对称轴的抛物线。


% ITEM: quadratic-exercise-16-02
\solutionitem{quadratic-exercise-16-02}{习题4.16（2）}
求以原点为焦点、\(x=-3\)为准线、离心率为\(1/2\)的曲线方程，并求中心与两条坐标方向上的半轴长度。
\par
\solutionanswer
设\(P=(x,y)\)。焦准线关系为\(\sqrt{x^2+y^2}=\frac12|x+3|\)，平方并配方得
\[
\frac34x^2+y^2-\frac32x-\frac94=0,
\qquad \frac{(x-1)^2}{4}+\frac{y^2}{3}=1.
\]
故曲线中心为\((1,0)\)；沿原\(x\)方向的半轴长为\(2\)，沿原\(y\)方向的半轴长为\(\sqrt3\)。


% ITEM: quadratic-exercise-16-03
\solutionitem{quadratic-exercise-16-03}{习题4.16（3）}
对\(x^2+y^2=4(x+1)^2\)，恢复其焦点、准线与离心率，化为标准方程，并求两条渐近线。
\par
\solutionanswer
原式正是焦点\(F=(0,0)\)、准线\(x=-1\)、\(d=1\)、离心率\(e=2\)的焦准线方程。整理并配方：
\[
x^2+y^2=4(x+1)^2
\iff -3\left(x+\frac43\right)^2+y^2=-\frac43
\iff \frac{(x+4/3)^2}{4/9}-\frac{y^2}{4/3}=1.
\]
双曲线中心为\((-4/3,0)\)。将原式给出的一组焦点、准线关于中心对称，即得另一焦点\((-8/3,0)\)及其对应准线\(x=-5/3\)。令右端常数项为零，渐近线满足\(-3X^2+Y^2=0\)，其中\(X=x+4/3,Y=y\)，故两条渐近线为\(y=\pm\sqrt3(x+4/3)\)。


% ITEM: quadratic-exercise-16-04
\solutionitem{quadratic-exercise-16-04}{习题4.16（4）}
由式\textup{(4.75)}推导椭圆与双曲线的半轴参数\(a,b\)与\(e,d\)的关系，再说明保持\(e\)不变、改变\(d\)时，为什么所得曲线相似。
\par

\par\smallskip\noindent 正文所引关系为：
\begin{equation*}
(1-e^2)\left(x-\frac{e^2d}{1-e^2}\right)^2+y^2
    =\frac{e^2d^2}{1-e^2} .
\end{equation*}
\solutionanswer
在焦点为原点、准线为\(x=-d\)的坐标中，焦准线方程化简后为
\[
(1-e^2)\left(x-\frac{e^2d}{1-e^2}\right)^2+y^2=\frac{e^2d^2}{1-e^2}.
\]
当\(0<e<1\)时，中心横坐标为\(e^2d/(1-e^2)\)，标准椭圆的半长轴沿\(x\)方向，且
\[
a=\frac{ed}{1-e^2},\qquad b=\frac{ed}{\sqrt{1-e^2}},\qquad b^2=a^2(1-e^2).
\]
当\(e>1\)时改写为
\[
(e^2-1)\left(x-\frac{e^2d}{1-e^2}\right)^2-y^2=\frac{e^2d^2}{e^2-1},
\]
中心仍为\(e^2d/(1-e^2)\)，横向实半轴及共轭半轴为
\[
a=\frac{ed}{e^2-1},\qquad b=\frac{ed}{\sqrt{e^2-1}},\qquad b^2=a^2(e^2-1).
\]
固定\(e\)后，\(a,b\)以及中心到焦点的距离都与\(d\)成正比；将\(d\)换成\(d'\)，曲线正是以焦点为中心按比例\(d'/d\)缩放所得，因此彼此相似。


% ITEM: quadratic-exercise-16-05
\solutionitem{quadratic-exercise-16-05}{习题4.16（5）}
对给定的\(e>1\)，选择一个可以实现它的圆锥半顶角\(\alpha\)，求相应截平面倾角。这样的圆锥是否唯一？
\par
\solutionanswer
给定\(e>1\)，可取圆锥半顶角\(\alpha=\arccos(1/e)\)，它满足\(0<\alpha<90^\circ\)。令截平面倾角\(\beta=90^\circ\)，则\(\sin\beta/\cos\alpha=e\)，确实得到该离心率的双曲线截线。圆锥并不唯一：任取\(\alpha\)满足\(\arccos(1/e)\le\alpha<90^\circ\)，再取\(\beta=\arcsin(e\cos\alpha)\)，都有\(0<\beta\le90^\circ\)且离心率为\(e\)，给出不同的可行圆锥半顶角。


% ITEM: quadratic-exercise-16-06
\solutionitem{quadratic-exercise-16-06}{习题4.16（6）}
求到\(F(1,1)\)的距离为到直线\(x+y=0\)距离的\(2\)倍的点的轨迹。选择以\(F\)为原点、垂直于准线的方向为一条坐标轴的坐标系，先写出新方程，再与原直角坐标中的方程比较。
\par
\solutionanswer
取以焦点\(F=(1,1)\)为原点、沿准线指向焦点的单位法向量\((1,1)/\sqrt2\)为第一坐标正方向，并取\((-1,1)/\sqrt2\)为第二坐标方向。记新坐标为\((u,v)\)，则
\[
u=\frac{x+y-2}{\sqrt2},\qquad v=\frac{y-x}{\sqrt2}.
\]
准线\(x+y=0\)在新坐标中为\(u=-\sqrt2\)，故\(d=\sqrt2\)、\(e=2\)。焦准线关系给出
\[
u^2+v^2=4(u+\sqrt2)^2.
\]
利用\(x+y=\sqrt2(u+\sqrt2)\)及\((x-1)^2+(y-1)^2=u^2+v^2\)，换回原坐标得
\[
(x-1)^2+(y-1)^2=2(x+y)^2,
\]
即\(x^2+y^2+4xy+2x+2y-2=0\)。此式也可直接由原坐标中的焦点距离平方等于\(4\)倍点到直线\(x+y=0\)的距离平方得到。
\par
